
\documentclass[11pt]{article}
\usepackage[margin=1in]{geometry}

% Core packages
\usepackage{amsmath,amssymb}
\usepackage{tikz-cd}
\usepackage{multicol}

% Paragraphs
\setlength{\parindent}{0pt}
\setlength{\parskip}{1\baselineskip}

\title{Benchmarking Statistical and Machine-Learning Models of Overall IEQ Satisfaction in Belgian Classrooms}
\author{Author Name}
\date{\today}

\begin{document}
\maketitle

\begin{abstract}
Provide a concise summary of the research problem, methods, main findings, and implications.
\end{abstract}

\section{Introduction}

\subsection{Context and Problem Definition}

Indoor environmental quality (IEQ) is a multidimensional concept that brings together thermal conditions, indoor air quality, lighting, and acoustics in order to describe how indoor environments are experienced by their occupants. In educational buildings, IEQ is not a purely technical concern. Classroom conditions shape occupants' comfort, satisfaction, and well-being, and they can also influence engagement and learning-related outcomes. For that reason, the question is no longer whether IEQ matters in classrooms, but how it should be modelled in a way that is both scientifically credible and practically useful for future building operation and occupant-centric control strategies \cite{heinzerling2013,wong2008,carton2023iaqvec}.

The literature has already shown that overall IEQ can be represented as more than the sum of isolated comfort domains. Early studies in offices and other non-residential settings demonstrated that overall acceptance or satisfaction can be estimated from combinations of thermal, air-quality, acoustic, and visual variables \cite{wong2008,frontczak2012}. At the same time, review work has emphasized that IEQ evaluation is methodologically heterogeneous: studies differ in their target variables, weighting strategies, measurement approaches, and validation protocols, which makes results difficult to compare directly \cite{heinzerling2013}. This fragmentation is especially problematic when predictive models are presented as generally applicable even though they may depend heavily on a specific dataset or a favorable data split.

A second challenge is that subjective satisfaction with indoor environments is not determined by environmental measurements alone. Recent review evidence on subjective indoor air quality in non-residential buildings shows that personal and contextual factors repeatedly influence reported satisfaction, and that predictive modelling in field settings is still relatively underdeveloped \cite{carton2026review}. This point is crucial for classroom research, where occupants differ in age, expectations, sensitivity, and exposure patterns, and where the same environmental conditions may be evaluated differently by different individuals. As a result, a purely generalized, sensor-only model may be insufficient if the goal is to predict overall IEQ satisfaction in a realistic and deployable way.

The present study addresses this problem through the open Belgian classroom dataset described by Carton et al. \cite{carton2026dataset}. The dataset contains 6{,}834 satisfaction assessments from 321 unique occupants across three educational settings: a secondary school, primary schools, and a university lecture room. For each assessment, repeated responses on thermal, IAQ, visual, acoustic, and overall IEQ satisfaction are available together with environmental and contextual information. This longitudinal structure creates an unusually strong basis for benchmarking generalized, occupant-informed, and personalized modelling strategies under conditions that approximate real deployment rather than idealized retrospective evaluation.

\subsection{Literature Review}

Research on overall IEQ and occupant satisfaction provides an important conceptual foundation for this study. Office-based studies have long shown that overall satisfaction depends on multiple IEQ domains at the same time rather than on thermal comfort alone \cite{wong2008,frontczak2012}. In parallel, broader review papers have highlighted the need for clearer IEQ assessment frameworks and more consistent interpretation of subjective and objective indicators \cite{heinzerling2013}. Together, these studies establish that multi-domain modelling is both meaningful and feasible, but they do not resolve how such modelling should be benchmarked in classroom environments with repeated occupant feedback.

School-focused studies reinforce the importance of the topic while also revealing the limits of the current evidence base. Sadick and Issa \cite{sadick2017} showed that teachers' satisfaction with ventilation and thermal comfort is closely related to dimensions of well-being, underlining the broader significance of IEQ in school buildings. In Belgian educational settings, Carton et al. reported frequent dissatisfaction with acoustic, IAQ, and thermal conditions in primary school classrooms and argued that classroom IEQ should be evaluated from the perspective of occupants' actual experience rather than technical compliance alone \cite{carton2023primary}. In a university classroom experiment, Carton et al. further showed that variations in thermal and IAQ conditions were associated with student satisfaction, engagement, and selected performance measures \cite{carton2023iaqvec}. These studies confirm that classroom IEQ is consequential, yet they are not designed as a rigorous head-to-head benchmark of predictive strategies for overall IEQ satisfaction.

The methodological literature also points toward the value of personalization. In office settings, Carton et al. \cite{carton2024mixed} found that mixed-effects models improved the prediction of individual thermal preferences relative to fixed-effects alternatives and suggested that only a limited number of prior feedback votes may be needed before personalized modelling becomes useful. Although this evidence comes from the thermal domain rather than from overall IEQ satisfaction, it provides a strong precedent for testing whether repeated occupant feedback adds predictive value in classroom applications. At the same time, the recent systematic review by Carton et al. \cite{carton2026review} emphasizes that subjective IEQ outcomes are often studied with inconsistent scales and analytical approaches, and that longitudinal repeated-feedback datasets remain rare. This strengthens the case for a transparent benchmark that explicitly compares generalized, occupant-informed, and personalized models.

Taken together, the literature suggests that the field does not lack proof that overall IEQ can be modelled. Instead, it lacks a disciplined benchmark that answers several unresolved questions at once: whether personalization improves prediction of overall IEQ satisfaction, whether hierarchical pipelines are preferable to direct prediction, and how strongly apparent model performance depends on the realism of the validation protocol. Existing studies usually address only one part of this problem. What is still missing is a leakage-safe comparison on an open classroom dataset that evaluates direct and hierarchical approaches under random, unseen-occupant, and unseen-classroom conditions while also examining transfer across educational contexts.

\subsection{Research Questions and Hypotheses}

In response to these gaps, this study asks how overall IEQ satisfaction in classrooms should be modelled, personalized, and validated when the goal is not merely high retrospective accuracy but credible performance under realistic deployment scenarios. The core research questions are formulated as follows:

\begin{enumerate}
    \item How accurately can overall IEQ satisfaction be predicted from a harmonized set of thermal, IAQ, visual, acoustic, contextual, and occupant-related variables?
    \item Do occupant-informed and personalized models outperform purely generalized models, and how many prior feedback votes are needed before personalization becomes beneficial?
    \item Does a hierarchical strategy improve predictive performance or interpretability compared with direct prediction of overall IEQ satisfaction?
    \item How strongly does predictive performance decrease under stricter evaluation settings, especially when models are tested on unseen occupants and unseen classrooms?
\end{enumerate}

The study also considers secondary questions related to transferability across age groups, classroom types, and, where feasible, external datasets. Based on the reviewed literature, four working hypotheses are proposed. First, models that combine environmental, contextual, and occupant-related variables are expected to outperform purely generalized sensor-based models because subjective IEQ satisfaction depends on more than physical measurements alone \cite{carton2026review}. Second, personalized models are expected to outperform generalized models once a sufficient number of prior feedback votes is available, consistent with evidence from mixed-effects comfort modelling \cite{carton2024mixed}. Third, hierarchical modelling is expected to improve interpretability and may improve predictive performance when domain-level satisfaction carries meaningful intermediate information, although deployable two-stage models should be evaluated separately from analytical upper-bound hierarchies to avoid leakage. Fourth, model performance is expected to decline under unseen-occupant and unseen-classroom validation compared with random splits, indicating that optimistic evaluation protocols can overstate real-world predictive performance.

\section{Methods}

\subsection{Research Design}

Describe the overall research design and justify why it is appropriate for the study.

\subsection{Sample / Participants}

Describe the research sample or participants, including relevant demographic or selection details.

\subsection{Data Collection Instruments}

Describe the tools, instruments, or procedures used to collect the data.

\subsection{Data Analysis Procedure}

Explain how the data were processed and analyzed.

\section{Results}

Present the data and statistical findings objectively, without interpretation.

\section{Discussion}

\subsection{Interpretation of Results and Answers to the Research Questions}

Interpret the findings and explain how they address the research questions.

\subsection{Comparison with Existing Literature}

Compare the results with prior studies and discuss similarities or differences.

\subsection{Limitations}

Identify the main limitations of the study.

\subsection{Recommendations for Future Research}

Suggest directions for future research based on the findings and limitations.

\section{Conclusion}

Provide a final summary of the main findings and their implications for practice or theory.

\begin{thebibliography}{99}

\bibitem{carton2023primary}
Q. Carton, F. Mennes, S. Vanden Broeck, V. Van Roy, J. Kolarik, and H. Breesch,
``Evaluation of indoor environmental quality and pupils' satisfaction in Flemish primary schools,''
in \emph{Proceedings of Healthy Buildings Europe 2023}, 2023.

\bibitem{carton2023iaqvec}
Q. Carton, S. De Coninck, J. Kolarik, and H. Breesch,
``Assessing the effect of a classroom IEQ on student satisfaction, engagement and performance,''
in \emph{E3S Web of Conferences}, vol. 396, 01052, 2023.

\bibitem{carton2024mixed}
Q. Carton, J. Kloppenborg Møller, M. Favero, D. Calì, J. Kolarik, and H. Breesch,
``Predicting individual thermal preferences in an office: Assessing the performance of mixed-effects models,''
\emph{Building and Environment}, vol. 261, 111751, 2024.

\bibitem{carton2026dataset}
Q. Carton, J. Kolarik, and H. Breesch,
``A dataset on occupant satisfaction with the indoor environmental quality in Belgian classrooms,''
\emph{Scientific Data}, vol. 13, 229, 2026.

\bibitem{carton2026review}
Q. Carton, D. Al-Assaad, J. Kolarik, and H. Breesch,
``A review of subjective indoor air quality assessment in non-residential buildings: Current trends and recommendations,''
\emph{Buildings}, vol. 16, 486, 2026.

\bibitem{frontczak2012}
M. Frontczak, S. Schiavon, J. Goins, E. Arens, H. Zhang, and P. Wargocki,
``Quantitative relationships between occupant satisfaction and satisfaction aspects of indoor environmental quality and building design,''
\emph{Indoor Air}, vol. 22, pp. 119--131, 2012.

\bibitem{heinzerling2013}
D. Heinzerling, S. Schiavon, T. Webster, and E. Arens,
``Indoor environmental quality assessment models: A literature review and a proposed weighting and classification scheme,''
\emph{Building and Environment}, 2013.

\bibitem{sadick2017}
A.-M. Sadick and M. H. Issa,
``Occupants' indoor environmental quality satisfaction factors as measures of school teachers' well-being,''
\emph{Building and Environment}, vol. 119, pp. 99--109, 2017.

\bibitem{wong2008}
L. T. Wong, K. W. Mui, and P. S. Hui,
``A multivariate-logistic model for acceptance of indoor environmental quality (IEQ) in offices,''
\emph{Building and Environment}, vol. 43, pp. 1--6, 2008.

\end{thebibliography}

\end{document}
